\chapter{Modelos e Instancias de Álgebras de Boole}

  Como ejemplos que cumplen los anteriores postulados o axiomas vamos a
  desarrollar unos cuántos. Para que este sistema de axiomas sea
  relativamente consistente (al sistema ZFC por ejemplo) bastará ver la
  independencia de unos axiomas de otros.

  \textbf{{[}Ejemplo 1{]}} El primero y más sencillo de ver es el
  álgebra de las partes de un conjunto. Dado un conjunto
  cualquiera\[U \neq \varnothing\],
  \[{\wp(U)} = {\{{X \mid {X \subseteq U}}\}}\], esto es,
  \[{\wp(U)} = {\{{X \mid {{\forall x}\left( {{({x \in X})}\Rightarrow{({x \in U})}} \right)}}\}}\],
  dónde se verifica
  que\[\left( {{\varnothing \in }\wp(U)} \right) \land \left( {{U \in }\wp(U)} \right)\].
  Ha\-remos\[\left( {B{: = }\wp(U)} \right),\left( {0{: = \varnothing}} \right)y\left( {1{: = U}} \right)\],
  como producto lógico pondremos la intersección de
  conjuntos\[\forall X,{Y \in }\wp(U){X \cdot Y}{: = {X \cap Y}}\], como
  suma lógica pondremos la unión de
  conjuntos\[\forall X,{Y \in }\wp(U){X + Y}{: = {X \cup Y}}\]. Las tres
  primeras (dobles) propiedades son di\-rectamente cumplidas por la
  estructura construida y la existencia del complementario es fá\-cil de
  ver. Sea
  \[\forall{X \in }\wp(U)\Rightarrow\exists Y{: = {U \smallsetminus X}}\]y a
  partir de ahí sabemos que\[{Y \in }\wp(U)\]puesto
  que\[\forall{x \in Y}{x \in {U \smallsetminus X}}\Rightarrow{x \in U}\]y
  en el caso que\[X = U\]tenemos que
  \[{{Y = {U \smallsetminus X}} = {U \smallsetminus U}} = \varnothing\]de
  forma que\[{\varnothing \in }\wp(U)\]por definición. Ahora sólo se trata
  de ver
  que\[{{{{Y \cdot X} = {Y \cap X}} = {{({U \smallsetminus X})} \cap X}} = \varnothing} = 0\]y
  que la propiedad dual a
  cumplir\[{{{{Y + X} = {Y \cup X}} = {{({U \smallsetminus X})} \cup X}} = U} = 1\].
  Ya tenemos
  que\[\forall{X \in }\wp(U)\exists{Y \in }\wp(U){Y \in \overline{X}}\]ha\-biendo
  tomado\[Y{: = {U \smallsetminus X}}\].

  \textbf{{[}Ejemplo 2{]}} Un ejemplo interesante fácil de construir es
  el álgebra de Boole de los números que son producto de los primeros
  núme\-ros primos (cantidad finita de ellos) y sus divisores.
  Consideramos el conjunto \[P_{n}{: = {\{{2,3,\ldots,p_{n}}\}}}\],
  con\-sideraremos el \[1\] booleano cómo
  \[1_{B}{: = {\prod\limits_{q \in P_{n}}q}}\]y el \[0\]
  cómo\[0_{B}{: = 1_{\mathbb{N}}}\]. Consideramos a
  \[B{: = {\{{{n \in \mathbb{N}} \mid {n \mid \left( {\prod P_{n}} \right)}}\}}}\],
  y las operaciones serán el mínimo común múltiplo como suma booleana y
  el máximo común divisor como producto booleano. El complemento de un
  ele\-mento resulta ser
  \[\forall{k \in B}{{\overline{k} = {1_{B}/k}} = {\prod\limits_{q \in {\{{{{p \in P_{n}} \mid p} \nmid k}\}}}q}}\].
  Éste es un modelo fácil de desarrollar para poner ejemplos.

  \textbf{{[}Ejemplo 3{]} }Partimos de un álgebra de Boole cualquiera y
  un elemento no \[0\] ni \[1\] cualquie\-ra tal
  que\[x \in \left( {B \smallsetminus {\{ 0,1\}}} \right)\]. Definimos
  ahora un
  conjunto\[B_{\leq x}{: = {\{{{{y \in B} \mid {x \cdot y}} = y}\}}}\]y\[B_{\geq x}{: = {\{{{{y \in B} \mid {x \cdot y}} = x}\}}}\].
  Las álgebras de Boole nuevas a considerar son
  \[\langle{B_{\leq x},{\{{{0 \equiv 0_{B}},{1 \equiv x}}\}},{\{{+_{B}, \cdot_{B}}\}}}\rangle\]
  y
  \[\langle{B_{\geq x},{\{{{0 \equiv x},{1 \equiv 1_{B}}}\}},{\{{+_{B}, \cdot_{B}}\}}}\rangle\].
  Consideramos el mismo producto booleano que en el conjunto inicial e
  idéntica\-mente con la suma booleana. Sólo varía el complemento, de la
  siguiente
  forma\[{y \in B_{\geq x}}\Rightarrow y{' = {\overline{y} + x}}\]y\[{y \in B_{\leq x}}\Rightarrow y{' = {\overline{y} \cdot x}}\].
  Es fácil comprobar la validez de esta definición de un álgebra de
  Boole, de manera más concreta, que las operaciones son internas y el
  complemento declarado es también interno y se comporta como
  complemento del nuevo álgebra, esto
  es,\[{y \in B_{\geq x}}\Rightarrow y{{' \cdot y} = {x \land y}}{{' + y} = 1}\]
  que\[{y \in B_{\leq x}}\Rightarrow y{{' \cdot y} = {0 \land y}}{{' + y} = x}\].

  \textbf{{[}Ejemplo 4{]}} El álgebra de las proposiciones. Éste es sin
  duda, el primer desarrollo que se hizo del álgebra de Boole, hecha por
  el propio George Boole en mitad del siglo XIX (edición 1851). Lo
  desarrolló como una ``Una investigación en las leyes del
  pensamiento''. Amigo suyo que lo ayudó a penetrar los ambientes
  académi\-cos es Augustus De Morgan. Desde Aristóteles (que funda por
  primera vez la ló\-gica como ciencia analítica) en el siglo IV a. C.
  no había habido ningún adelanto sustancial en la lógica. Kant (medio
  siglo antes de Boole) había considerado que la lógica era un cuerpo de
  doctrina cerrado y completo (esto es, no había nada más que decir que
  lo que ya había desarrollado y escrito Aristóteles en sus ``Tra\-tados
  de Lógica'' u ``Órganon'', hacía ya 2.300 años). Aunque la verdadera
  revolu\-ción se da algunos años más tarde con Frege, el lógico más
  importante desde Aristóteles. Si doy estos datos sobre la historia de
  la lógica que todos asociaréis más a la filosofía, que parece queda
  muy lejos del propósito de unos apuntes de matemáticas discretas que
  cubran de la forma más amplia posible los Fundamentos de Electró\-nica
  en su parte Digital, es porque no queda tan lejos. La idea de hacer un
  len\-guaje dónde el razonamiento siguiera unas pautas claras de forma
  que siempre quedara todo tan cierto como en las matemáticas era ya
  antiguo. Aristóteles ya advertía de una cierta in-definición
  insuperable de los términos más importantes de la filosofía (en
  realidad de casi todos los conceptos de la vida ordinaria):
  ``exis\-ten conceptos o ideas que corresponden con la realidad que no
  se usan de forma equívoca -- esto es, su uso no es equívoco, este
  mismo concepto, palabra o idea que hablamos no se refiere a realidades
  distintas y diferenciadas, de forma que nos llevan a confusión -- pero
  tampoco de forma unívoca -- como las definiciones desde axiomas en un
  lenguaje formal matemático, así tenemos que existen con\-ceptos
  análogos'' (es una glosa de palabras de Aristóteles). En la
  Modernidad, dado que el concepto de analogía lleva aparejado un
  tratamiento difícil que no lleva fácilmente a certeza, se intentan
  buscar criterios de certeza absoluta y unas definiciones que
  aparentemente son unívocas y se tratan como tales. Es significativo el
  nombre (y la estructura interna) de una importante obra de Spino\-za:
  ``Ética demostrada según el orden geométrico''. Será Leibniz quién
  escriba ya cumplidamente sobre la necesidad de establecer un léxico
  completamente unívoco (una tarea mastodonte, o mejor, imposible) y un
  ``cálculo'' del pensamiento, de forma que ``una cuestión como la
  existencia de Dios pueda ser resuelto mediante la resolución de unas
  ecuaciones de pensamiento'' (de nuevo es una glosa). Se empezaba a
  buscar con ansiedad una mecanización del pensamiento. Esta idea fue
  muy fructífera, dando un primer paso hacia ella George Boole que hace
  un ál\-gebra de las proposiciones. Esta álgebra no era más amplia que
  la de Aristóteles, pero permitía el cálculo al modo matemático. De
  aquí a la llegada de Frege, ya, Charles Babbage diseña y realiza (sin
  éxito debido al trabajo de mecanizado ex\-cesivamente minucioso que
  requería el diseño) una computadora universal me\-cánica (mediante
  engranajes) prácticamente similar al modelo de Von Neumann. La condesa
  de Lovelace (Ada) es el matemático que hace los primeros progra\-mas
  en lenguaje ensamblador de la máquina de Babbage. La primera
  programa\-dora de la historia. Después de Frege siguen los desarrollos
  con gente como Ber\-trand Russell, David Hilbert y otros, hasta los
  increíbles resultados de Gödel que ponen punto final a muchas de las
  pretensiones de mecanización del pensamiento, pero que son ya base de
  la computación moderna, siendo los trabajos definitivos los de Alan
  Turing. Como veis el camino recorrido es largo y complicado, siendo el
  momento crucial para el arranque de la ingeniería digital los trabajos
  de George Boole. No he mencionado el papel de las máquinas de cifrado
  y descifrado de mensajes en la Gran Guerra y la II Guerra Mundial (en
  las que intervinieron muchos de las mentes antes mencionadas).

  De manera un tanto informal podemos ver una proposición (una frase que
  afirma o niega una propiedad de un objeto, una relación entre objetos
  o la existencia del mismo, una frase que ha de ser o verdadero,
  \[1_{B}{{: = V} \equiv \mathbf{\mathit{true}}}\], a falso,
  \[0_{B}{{: = F} \equiv \mathbf{\mathit{false}}}\]) o conjunto de
  proposiciones pueden ser operadas mediante la conjunción `y', A `y' B
  es verdadero si A es verdadero y B es verdadero a la vez y falso en
  cualquier otro caso. La disyunción sería la `o', siendo A `o' B
  verdadero con que A sea verdadero o lo sea B, siendo falso sólo cuando
  A es falso y B es falso a la vez. La notación más habitual es
  \[{{+ {: = \vee}} \equiv \text{or}} \equiv {\mid \mid}\]y\[{{\cdot {: = \land}} \equiv \text{and}} \equiv {\&\&}\].
  Para el `no' (negación) tenemos
  que\[{{\overline{\phantom{A}}{: = {\neg\phantom{A}}}} \equiv \text{not}}{\phantom{A} \equiv {/\phantom{A}}}\].
  El conjunto de Boole es el conjunto de proposiciones de la que
  partamos.

  De manera más formal se considera un elemento del álgebra de Boole de
  la lógica a las clases de equivalencia de las proposiciones
  equivalentes lógicamente (en su valor de verdad o falsedad) entre sí.

  \textbf{{[}Ejemplo 5{]}} El álgebra de conmutación. Este es el álgebra
  de Boole más sencillo que hay. \[B{{: = B_{2}} \equiv {\{{0,1}\}}}\].
  Las operaciones las concretaremos en tablas:

  \setcounter{MaxMatrixCols}{20}
  \setlength{\arraycolsep}{8pt} % Aumentado a peticion del usuario
  
enewcommand{rraystretch}{2}
  \[\begin{bmatrix}
   + & 0 & 1 \\
  0 & 0 & 1 \\
  1 & 1 & 1
  \end{bmatrix}\begin{bmatrix}
   \cdot & 0 & 1 \\
  0 & 0 & 0 \\
  1 & 0 & 1
  \end{bmatrix}\begin{bmatrix}
  \overline{} & 0 & 1 \\
   & 1 & 0
  \end{bmatrix}\]

  y podréis comprobar fácilmente que se cumplen todos los postulados de
  Huntington. Ésta será usada frecuentemente durante el curso. Esta
  álgebra está contenido en todo álgebra de Boole.

  \textbf{{[}Ejemplo 6{]}} El álgebra de Boole de 4 elementos. Este es
  el álgebra de Boole generada por un conjunto de 2 elementos. Es
  singular en el sentido que sólo tiene 3 niveles, el más bajo
  \[\{{0{{: = {\{\}}} \equiv \varnothing}}\}\], el intermedio
  \[\{{{\{\alpha\}},{\{\beta\}}}\}\], y el superior
  \[\{{1{: = {\{{\alpha,\beta}\}}}}\}\].
  \[B{{: = B_{4}} \equiv {\{{0,a,b,1}\}}}\]. Las operaciones las
  concretaremos en tablas:

  \[\begin{bmatrix}
   + & 0 & a & b & 1 \\
  0 & 0 & a & b & 1 \\
  a & a & a & 1 & 1 \\
  b & b & 1 & b & 1 \\
  1 & 1 & 1 & 1 & 1
  \end{bmatrix}\begin{bmatrix}
   + & 0 & a & b & 1 \\
  0 & 0 & 0 & 0 & 0 \\
  a & 0 & a & 0 & a \\
  b & 0 & 0 & b & b \\
  1 & 0 & a & b & 1
  \end{bmatrix}\begin{bmatrix}
  \overline{} & 0 & a & b & 1 \\
   & 1 & b & a & 0
  \end{bmatrix}\]

  y podréis comprobar fácilmente que se cumplen todos los postulados de
  Huntington si cam\-biáis \[a\] por \[\{\alpha\}\], \[b\] por
  \[\{\beta\}\], \[1\] por \[\{{{\{\alpha\}},{\{\beta\}}}\}\] y \[0\]
  por el conjunto vacío \[\varnothing\].

  \textbf{{[}Ejemplo 7{]}} El álgebra de Boole de 8 elementos. Este es
  el álgebra de Boole generada por un conjunto de 3 elementos. Es
  singular en el sentido que sólo tiene 4 niveles, el más bajo
  \[\{ 0\}\], el de átomos \[\{{a,b,c}\}\], el de hiperátomos
  \[\{{A,B,C}\}\] y el superior \[\{ 1\}\]. Los niveles de átomos y de
  hiperátomos son especialmente importantes, siendo esta álge\-bra de
  Boole, la más pequeña que los diferencia. Sería:

  \[B{{: = B_{8}} \equiv {\{{0,a,b,c,A,C,B,1}\}}}\].

  Las operaciones las concretaremos en tablas:

  \setcounter{MaxMatrixCols}{20}
  \setlength{\arraycolsep}{8pt} % Aumentado a peticion del usuario
  
enewcommand{rraystretch}{2}
  \[\begin{bmatrix}
   + & 0 & a & b & c & A & C & B & 1 \\
  0 & 0 & a & b & c & A & C & B & 1 \\
  a & a & a & A & C & A & C & 1 & 1 \\
  b & b & A & b & B & A & 1 & B & 1 \\
  c & c & C & B & c & 1 & C & B & 1 \\
  A & A & A & A & 1 & A & 1 & 1 & 1 \\
  C & C & C & 1 & C & 1 & C & 1 & 1 \\
  B & B & 1 & B & B & 1 & 1 & B & 1 \\
  1 & 1 & 1 & 1 & 1 & 1 & 1 & 1 & 1
  \end{bmatrix}\begin{bmatrix}
   \cdot & 0 & a & b & c & A & C & B & 1 \\
  0 & 0 & 0 & 0 & 0 & 0 & 0 & 0 & 0 \\
  a & 0 & a & 0 & 0 & a & a & 0 & a \\
  b & 0 & 0 & b & 0 & b & 0 & b & b \\
  c & 0 & 0 & 0 & c & 0 & c & c & c \\
  A & 0 & a & b & 0 & A & a & b & A \\
  C & 0 & a & 0 & c & a & C & c & C \\
  B & 0 & 0 & b & c & b & c & B & B \\
  1 & 0 & a & b & c & A & C & B & 1
  \end{bmatrix}\]

  \[\begin{bmatrix}
  \neg & 0 & a & b & c & A & C & B & 1 \\
   & 1 & B & C & A & c & b & a & 0
  \end{bmatrix}\begin{Bmatrix}
  {{B_{8}\rightarrow\wp}{\{{\alpha,\beta,\gamma}\}}} & {0\rightarrow\varnothing} \\
  {a\rightarrow{\{\alpha\}}} & {b\rightarrow{\{\beta\}}} \\
  {c\rightarrow{\{\gamma\}}} & {A\rightarrow{\{{\alpha,\beta}\}}} \\
  {C\rightarrow{\{{\gamma,\alpha}\}}} & {B\rightarrow{\{{\beta,\gamma}\}}} \\
  {1\rightarrow{\{{\alpha,\beta,\gamma}\}}} & {{x \cdot y}\rightarrow{x \cap y}} \\
  {{x + y}\rightarrow{x \cup y}} & {\overline{x}\rightarrow{{\{{\alpha,\beta,\gamma}\}} \smallsetminus x}}
  \end{Bmatrix}\]

  y podréis comprobar fácilmente que se cumplen todos los postulados de
  Huntington si te\-néis en cuenta los cambios aconsejados en el cuadro
  entre llaves, dónde las flechas quieren decir ``substituir por''.

  \textbf{{[}Ejemplo 8{]}} El álgebra de Boole de 16 elementos. Este es
  el álgebra de Boole generada por un conjunto de 4 elementos. Es ya un
  álgebra de Boole completamente regular. Tiene 5 niveles, el más bajo
  el \[\{ 0\}\], el de átomos \[\{{\alpha,\beta,\gamma,\delta}\}\], el
  de hiperátomos \[\{{A,B,\Gamma,\Delta}\}\], el intermedio
  \[\{{a,b,c,d,e,f}\}\]y finalmente el nivel superior con el \[\{ 1\}\].
  Sería:

  \[B{{: = B_{16}} \equiv {\{{0,\alpha,\beta,\gamma,\delta,a,b,c,d,e,f,A,B,\Gamma,\Delta,1}\}}}\].

  Las operaciones las concretaremos en tablas:

  \setcounter{MaxMatrixCols}{20}
  \setlength{\arraycolsep}{8pt} % Aumentado a peticion del usuario
  
enewcommand{rraystretch}{2}
  \[\begin{bmatrix}
   + & 0 & \alpha & \beta & \gamma & \delta & a & b & c & d & e & f & A & B & \Gamma & \Delta & 1 \\
  0 & 0 & \alpha & \beta & \gamma & \delta & a & b & c & d & e & f & A & B & \Gamma & \Delta & 1 \\
  \alpha & \alpha & \alpha & a & b & c & a & b & c & A & B & \Gamma & A & B & \Gamma & 1 & 1 \\
  \beta & \beta & a & \beta & d & e & a & A & B & d & e & \Delta & A & B & 1 & \Delta & 1 \\
  \gamma & \gamma & b & d & \gamma & f & A & b & B & \Gamma & \Delta & f & A & 1 & \Gamma & \Delta & 1 \\
  \delta & \delta & c & e & f & \delta & B & \Gamma & c & \Delta & e & f & 1 & B & \Gamma & \Delta & 1 \\
  a & a & a & a & A & B & a & A & B & A & \Delta & 1 & A & B & 1 & 1 & 1 \\
  b & b & b & A & b & \Gamma & A & b & B & A & 1 & \Gamma & A & 1 & \Gamma & 1 & 1 \\
  c & c & c & B & \Gamma & c & B & \Gamma & c & 1 & B & \Gamma & 1 & B & 1 & \Delta & 1 \\
  d & d & A & d & d & \Delta & A & A & 1 & d & \Delta & \Delta & A & 1 & 1 & \Delta & 1 \\
  e & e & B & e & \Delta & e & B & 1 & B & \Delta & e & \Delta & 1 & B & 1 & \Delta & 1 \\
  f & f & \Gamma & \Delta & f & f & 1 & \Gamma & \Gamma & \Delta & \Delta & f & 1 & 1 & \Gamma & \Delta & 1 \\
  A & A & A & A & A & 1 & A & A & 1 & A & 1 & 1 & A & 1 & 1 & 1 & 1 \\
  B & B & B & B & 1 & B & B & 1 & B & 1 & B & 1 & 1 & B & 1 & 1 & 1 \\
  \Gamma & \Gamma & \Gamma & 1 & \Gamma & \Gamma & 1 & \Gamma & \Gamma & 1 & 1 & \Gamma & 1 & 1 & \Gamma & 1 & 1 \\
  \Delta & \Delta & 1 & \Delta & \Delta & \Delta & 1 & 1 & 1 & \Delta & \Delta & \Delta & 1 & 1 & 1 & \Delta & 1 \\
  1 & 1 & 1 & 1 & 1 & 1 & 1 & 1 & 1 & 1 & 1 & 1 & 1 & 1 & 1 & 1 & 1
  \end{bmatrix}\]

\[\begin{bmatrix}
\neg & 0 & \alpha & \beta & \gamma & \delta & a & b & c & d & e & f & A & B & \Gamma & \Delta & 1 \\
 & 1 & \Delta & \Gamma & B & A & f & e & d & c & b & a & \delta & \gamma & \beta & \alpha & 0
\end{bmatrix}\]

\[\begin{bmatrix}
 \cdot & 0 & \alpha & \beta & \gamma & \delta & a & b & c & d & e & f & A & B & \Gamma & \Delta & 1 \\
0 & 0 & 0 & 0 & 0 & 0 & 0 & 0 & 0 & 0 & 0 & 0 & 0 & 0 & 0 & 0 & 0 \\
\alpha & \alpha & \alpha & 0 & 0 & 0 & \alpha & \alpha & \alpha & 0 & 0 & 0 & \alpha & \alpha & \alpha & 0 & \alpha \\
\beta & \beta & 0 & \beta & 0 & 0 & \beta & 0 & 0 & \beta & \beta & 0 & \beta & \beta & 0 & \beta & \beta \\
\gamma & \gamma & 0 & 0 & \gamma & 0 & 0 & \gamma & 0 & \gamma & 0 & \gamma & \gamma & 0 & \gamma & \gamma & \gamma \\
\delta & \delta & 0 & 0 & 0 & \delta & 0 & 0 & \delta & 0 & \delta & \delta & 0 & \delta & \delta & \delta & \delta \\
a & a & \alpha & \beta & 0 & 0 & a & \alpha & \alpha & \beta & \beta & 0 & a & a & \alpha & \beta & a \\
b & b & \alpha & 0 & \gamma & 0 & \alpha & b & \alpha & \gamma & 0 & \gamma & a & a & \alpha & \beta & b \\
c & c & \alpha & 0 & 0 & \delta & \alpha & \alpha & c & 0 & \delta & \delta & \alpha & c & c & \delta & c \\
d & d & 0 & \beta & \gamma & 0 & \beta & \gamma & 0 & d & \beta & \gamma & d & \beta & \gamma & d & d \\
e & e & 0 & \beta & 0 & \delta & \beta & 0 & \delta & \beta & e & \delta & \beta & e & \delta & e & e \\
f & f & 0 & 0 & \gamma & \delta & 0 & \gamma & \delta & \gamma & \delta & f & \gamma & \delta & f & f & f \\
A & A & \alpha & \beta & \gamma & 0 & a & b & \alpha & d & e & \gamma & A & a & b & d & A \\
B & B & \alpha & \beta & 0 & \delta & a & \alpha & c & \beta & e & \delta & a & B & c & d & B \\
\Gamma & \Gamma & \alpha & 0 & \gamma & \delta & \alpha & b & c & \gamma & \delta & f & b & c & \Gamma & f & \Gamma \\
\Delta & \Delta & 0 & \beta & \gamma & \delta & \beta & \gamma & \delta & d & e & f & d & e & f & \Delta & \Delta \\
1 & 1 & \alpha & \beta & \gamma & \delta & a & b & c & d & e & f & A & B & \Gamma & \Delta & 1
\end{bmatrix}\]

  y podréis comprobar fácilmente que se cumplen todos los postulados de
  Huntington, con solo tener en cuenta que todos los elementos se pueden
  poner en función de \[\alpha\beta\gamma\delta\]y sumas de ellos. Las
  sumas de dos de los anteriores elementos son \[abcdef\]y las sumas de
  tres de ellos son \[AB\Gamma\Delta\].

  \textbf{{[}Ejemplo 9{]}} El álgebra de Boole de los conjuntos que se
  pueden expresar como \textbf{unión des\-junta finita de subintervalos
  genéricos de \[\lbrack 0,1\rbrack \cap \mathbb{Q}\]. Definimos por
  conveniencia
  \[\mathbf{\mathrm{I}}_{\mathbb{Q}} ≝ \left\lbrack {0,1} \right\rbrack_{\mathbb{Q}}\]}.
  Para esto haremos abs\-tracción de cualquier conjunto finito de puntos
  de \textbf{\[\mathbf{\mathrm{I}}_{\mathbb{Q}}\]}, esto es,
  consideraremos que dos conjuntos son iguales si su diferencia
  simétrica (la unión de las diferencias, los elementos que no son
  comunes de ambos conjuntos) es vacía o es un conjunto finito de
  puntos. Esta álgebra de Boole tiene un cardinal infinito numerable
  (como el cardinal de los números naturales). Lo más importante es que
  no puede desarrollarse de manera semejante a como desarrolla\-mos el
  álgebra de las partes de un conjunto. Lo formalizaremos del siguiente
  modo:

  \begin{enumerate}
  \def\labelenumii{\arabic{enumii}.}
  \setcounter{enumii}{1}
  \item
    \[a,{b \in \mathbf{\mathrm{I}}_{\mathbb{Q}}}{a < b}\Rightarrow\left\lbrack {a,b} \right\rbrack_{\mathbb{Q}} ≝ {\left\lbrack {a,b} \right\rbrack \cap \mathbb{Q}} ≝ \left\{ {{x \in \mathbf{\mathrm{I}}_{\mathbb{Q}}} \mid {{a \leq x} \leq b}} \right\}\]

    \begin{enumerate}
    \def\labelenumiii{\arabic{enumiii}.}
    \item
      Si escribimos \[\left\lbrack {a,b} \right\rbrack_{\mathbb{Q}}\]
      entonces \[{a < {b \land a}} \neq b\].
    \item
      Sea
      \[\mathbf{II}_{\mathbb{Q}} ≝ \left\{ {\left\lbrack {a,b} \right\rbrack_{\mathbb{Q}} \mid {a,{{{b \in {I_{\mathbb{Q}} \land a}} < {b \land a}} \neq b}}} \right\}\].
    \item
      Sea
      \[\mathbf{\mathrm{III}}_{\mathbb{Q}} ≝ {\left\{ {{A \in {\wp\left( I_{\mathbb{Q}} \right)}} \mid {{A = \mathbf{\cup}_{\lambda \in \Lambda}}I_{\lambda}\forall{\lambda \in \Lambda}{I_{\lambda} \in {\mathbf{\mathrm{II}}_{\mathbb{Q}}{{\#\left( \Lambda \right)} \in \widetilde{\mathbb{N}}}}}}} \right\} \cup \left\{ \varnothing \right\}}\]
      .
    \item
      Sea
      \[\mathit{Fin}\left( I_{\mathbb{Q}} \right) ≝ \left\{ {{A \in \wp}\left( I_{\mathbb{Q}} \right) \mid \#{(A) \in \widetilde{\mathbb{N}}}} \right\}\].
    \item
      \[A,B{\in}{\wp\left( \mathbf{\mathrm{I}}_{\mathbb{Q}} \right)}{A \approx B} ≝ {\#{\left( {A \symdiff B} \right) \in \widetilde{\mathbb{N}}}}\].
      Esta relación es de equivalencia.

      \begin{enumerate}
      \def\labelenumiv{\arabic{enumiv}.}
      \item
        Reflexiva
        \[\#{\left( {A \symdiff A} \right) = \#}{{(\varnothing) = 0} \in \widetilde{\mathbb{N}}}\].
        Luego \[A \approx A\].
      \item
        Simétrica
        \[A \symdiff {B = B} \symdiff A.\Rightarrow.A \approx B\Leftrightarrow B \approx A\].
      \item
        Transitiva
        \[A \approx {B \land B} \approx C\Rightarrow A \approx C\].

        \begin{itemize}
        \tightlist
        \item
          \[\#{{\left( {A \symdiff B} \right) = n_{1}} \in {\widetilde{\mathbb{N}} \land \#}}{{\left( {B \symdiff C} \right) = n_{2}} \in \widetilde{\mathbb{N}}}.\Rightarrow.\#{{\left( {A \symdiff C} \right) \leq {n_{1} + n_{2}}} \in \widetilde{\mathbb{N}}}\].
          Y queda de\-mostrada la propiedad transitiva.
        \end{itemize}
      \end{enumerate}
    \item
      A partir de aquí hablaremos de \[⟦A⟧\]para hablar de la clase de
      equivalencia de \[A \in \mathbf{\mathrm{III}}_{\mathbb{Q}}\]bajo
      la relación de equivalencia \[\approx\].
    \item
      A partir de aquí hablaremos de nuestro conjunto
      \[\mathbf{\mathrm{I}}_{\mathbb{Q}}^{\mathbf{\mathrm{GEN}}}(0,1) ≝ \left\{ {{⟦A⟧} \mid {A \in \mathbf{\mathrm{III}}_{\mathbb{Q}}}} \right\}\]
    \item
      Nuestro conjunto de Boole será
      \[B ≝ {\mathbf{\mathrm{I}}_{\mathbb{Q}}^{\mathbf{\mathrm{GEN}}}(0,1)}\].
    \item
      El \[0 ≝ {⟦\varnothing ⟧}\].
    \item
      El \[1 ≝ {⟦\mathbf{\mathrm{I}}_{\mathbb{Q}}⟧}\].
    \item
      Ahora veremos unas operaciones muy cercanas a la unión, la
      intersección y el com\-plemento, que realmente nos dan un álgebra
      de Boole sobre
      \[\mathbf{\mathrm{I}}_{\mathbb{Q}}^{\mathbf{\mathrm{GEN}}}\]:

      \[\begin{matrix}
      {{⟦A⟧},{{⟦B⟧} \in \mathbf{\mathrm{I}}_{\mathbb{Q}}^{\mathbf{\mathrm{GEN}}}}} \\
      {{{⟦A⟧} + {⟦B⟧}} ≝ {⟦{A \cup B}⟧}} \\
      {{{⟦A⟧} \cdot {⟦B⟧}} ≝ {⟦{A \cap B}⟧}} \\
      {\overline{⟦A⟧} ≝ {⟦{\lbrack 0,1\rbrack_{\mathbb{Q}} \smallsetminus A}⟧}}
      \end{matrix}\]
    \item
      Convenio de
      notación:\[{⟦{a,b}⟧} ≝ {⟦\left\lbrack {a,b} \right\rbrack_{\mathbb{Q}}⟧}\].
      Estos conjuntos serán nuestros subintervalos genéricos del
      intervalo genérico unidad.
    \item
      Sea una sucesión finita de un número par \[2 \cdot n\] de
      elementos de \[\lbrack 0,1\rbrack_{\mathbb{Q}}\], estricta\-mente
      creciente
      \[{{{{{{{0_{\mathbb{Q}} \leq a_{1}} < b_{1}} < a_{2}} < b_{2}} < \ldots} < a_{n}} < b_{n}} \leq 1_{\mathbb{Q}}\]
      dispuestos como

      \[⟦{a_{1},b_{1},a_{2},b_{2},\ldots,a_{n},b_{n}}⟧\]definirán los
      elementos de
      \[\mathbf{\mathrm{I}}_{\mathbb{Q}}^{\mathbf{\mathrm{GEN}}}\],
      aparte de
      \[{⟦⟧} ≝ {{{{⟦\varnothing ⟧} = {⟦{\{ 0\}}⟧}} = {⟦{\{ 1\}}⟧}} = 0_{\mathbf{\mathrm{I}}_{B}^{\mathbf{\mathrm{GEN}}}}}\].
    \item
      Si escribimos \[⟦{a_{1},b_{1},a_{2},b_{2},\ldots,a_{n},b_{n}}⟧\],
      significamos ya (suponemos que es un hecho que)
      \[{{{{{{{0_{\mathbb{Q}} \leq a_{1}} < b_{1}} < a_{2}} < b_{2}} < \ldots} < a_{n}} < b_{n}} \leq 1_{\mathbb{Q}}\].
    \item
      Ahora ya definimos (notación):

      \[\begin{matrix}
      {{⟦{a_{1},b_{1},a_{2},b_{2},\ldots,a_{n},b_{n}}⟧} ≝ {{⟦{\{{{x \in {\lbrack 0,1\rbrack}_{\mathbb{Q}}} \mid \exists{{{1 \leq k} \leq n} \in \mathbb{N}}{x \in {\lbrack{a_{k},b_{k}}\rbrack}_{\mathbb{Q}}}}\}}⟧} \equiv}} \\
      {\equiv {⟦{\mathbf{\cup}_{k = 1}^{n}{\lbrack{a_{k},b_{k}}\rbrack}}⟧}}
      \end{matrix}\].
    \item
      Ahora ya tenemos el conjunto de Boole que buscábamos:

      \[\left\lbrack {⟦0,1⟧} \right\rbrack_{\mathbf{\mathrm{I}}} ≝ {\left\{ {{⟦{a_{1,}b_{1,}\ldots,a_{n},b_{n}}⟧} \mid {\exists{n \in \mathbb{N}}{{{{{{0_{\mathbb{Q}} \leq a_{1}} < b_{1}} < \ldots} < a_{n}} < b_{n}} \leq 1_{\mathbb{Q}}}}} \right\} \cup \left\{ {⟦⟧} \right\}}\].
    \end{enumerate}
  \end{enumerate}

  Que las uniones, complementos e intersecciones de intervalos genéricos
  finitos siguen siendo intervalos genéricos finitos es claro desde el
  principio. Sin embargo voy a exponer la cabalística, hacer las cuentas
  vamos, para que no quede lugar a dudas. Con toda esta comprobación (o
  re-definición) de que \[B\] es cerrado bajo las distintas operaciones
  es laborioso, un tanto enojoso.

  La operación de complemento queda de la siguiente manera, y aunque aún
  no podemos comprobar aún su corrección, si queda claro que es un
  operación unaria interna:

  \[
A^{I} := \begin{cases}
  \llbracket 0, a_1, b_1, a_2, \ldots, b_{n-1}, a_n, b_n, 1 \rrbracket & \text{si } a_1 \neq 0 \land b_n \neq 1 \\
  \llbracket b_1, a_2, \ldots, b_{n-1}, a_n, b_n, 1 \rrbracket & \text{si } a_1 = 0 \land b_n \neq 1 \\
  \llbracket b_1, a_2, \ldots, b_{n-1}, a_n \rrbracket & \text{si } a_1 = 0 \land b_n = 1 \\
  \llbracket 0, a_1, b_1, a_2, \ldots, b_{n-1}, a_n \rrbracket & \text{si } a_1 \neq 0 \land b_n = 1 \\
  \llbracket \varnothing \rrbracket & \text{si } A = \llbracket 0, 1 \rrbracket \\
  \llbracket 0, 1 \rrbracket & \text{si } A = \llbracket \varnothing \rrbracket
\end{cases}
\]

De dónde obtenemos $\forall A \in B, \exists A^{I} \in B$.

Tenemos que $0 \in B, 0 := \llbracket \varnothing \rrbracket \equiv \llbracket \rrbracket$ y $1 \in B, 1 := \llbracket 0,1 \rrbracket$, y $0^{I} = 1 \land 1^{I} = 0$. Además observamos con claridad que $\forall A \in B, \exists A^{I} \in B$ tal que $A + A^{I} = \llbracket \mathbf{\mathrm{I}}_{\mathbb{Q}} \rrbracket = 1$ y $A \cdot A^{I} = \llbracket \rrbracket = 0$. Además de $\forall A \in B, (A^{I})^{I} = A$. Así nos queda $A^{I} \equiv \overline{A}$ si se verifican los demás axiomas.

La suma quedará de la siguiente forma:

\[ \forall A, B \in \left\lbrack \llbracket 0,1 \rrbracket \right\rbrack_{\mathbf{\mathrm{I}}} \exists (A,B) \subset (A \times B), A + B \triangleq \llbracket A \cup B \rrbracket \]

El producto seguirá un camino par:

\[ \forall A, B \in \left\lbrack \llbracket 0,1 \rrbracket \right\rbrack_{\mathbf{\mathrm{I}}} \exists (A,B) \subset (A \times B), A \cdot B \triangleq \llbracket A \cap B \rrbracket \]

Sólo queda ver que efectivamente las operaciones son internas:

Prueba:

\begin{enumerate}
\item Ahora vamos a desarrollar la suma de forma recurrente:

\[ B \in \left\lbrack \llbracket 0,1 \rrbracket \right\rbrack_{\mathbf{\mathrm{I}}}, B = \llbracket a_{1}^{B},b_{1}^{B},\ldots,a_{n}^{B},b_{n}^{B} \rrbracket \]
\[ A = \llbracket a_{1}^{A},b_{1}^{A},\ldots,a_{m}^{A},b_{m}^{A} \rrbracket \]

Comenzaremos por $m = 0$ y algunos casos especiales:

\[
B + A := \begin{cases}
  \llbracket \varnothing \rrbracket & \text{si } A = \llbracket \varnothing \rrbracket \land B = \llbracket \varnothing \rrbracket \\
  A & \text{si } B = \llbracket \varnothing \rrbracket \\
  B & \text{si } A = \llbracket \varnothing \rrbracket \\
  1 & \text{si } \exists A \in A, \exists B \in B : \overline{B} \subseteq A \lor \overline{A} \subseteq B \\
  A & \text{si } \exists A \in A, \exists B \in B : B \subseteq A \\
  B & \text{si } \exists A \in A, \exists B \in B : A \subseteq B
\end{cases}
\]

Para el caso general de $m = 1$:

\begin{align*}
B + A &:= \begin{cases}
  \llbracket a_1^A, b_1^A, a_1^B, b_1^B, \ldots, a_n^B, b_n^B \rrbracket & \text{si } b_1^A < a_1^B \\
  \llbracket a_1^A, b_1^B, a_2^B, b_2^B, \ldots, a_n^B, b_n^B \rrbracket & \text{si } b_1^A \ge a_1^B \land b_1^A \le b_1^B \\
  \llbracket a_1^B, b_1^B, \ldots, a_n^B, b_n^B, a_1^A, b_1^A \rrbracket & \text{si } a_1^A > b_n^B \\
  \llbracket a_1^B, b_1^B, \ldots, a_{n-1}^B, b_{n-1}^B, a_n^B, b_1^A \rrbracket & \text{si } a_1^A \le b_n^B \land a_1^A \ge a_n^B \land b_1^A > b_n^B \\
  \llbracket a_1^A, b_1^A, a_k^B, b_k^B, \ldots, a_n^B, b_n^B \rrbracket & \text{si } a_1^A \le a_1^B \land \exists k<n : b_1^A > b_{k-1}^B \land b_1^A < a_k^B \\
  \llbracket a_1^A, b_k^B, a_{k+1}^B, b_{k+1}^B, \ldots, a_n^B, b_n^B \rrbracket & \text{si } a_1^A \le a_1^B \land \exists k<n : b_1^A \ge a_k^B \land b_1^A \le b_k^B
\end{cases} \\[1em]
&\phantom{:=} \begin{cases}
  \llbracket a_1^B, b_1^B, \ldots, b_{k-1}^B, a_1^A, b_1^A, a_k^B, b_k^B, \ldots \rrbracket & \text{si } \exists k<n : a_1^A > b_{k-1}^B \land a_1^A \le a_k^B \land b_1^A < a_k^B \\
  \llbracket a_1^B, b_1^B, \ldots, a_l^A, b_k^B, a_{k+1}^B, \dots \rrbracket & \text{si } \exists l<k<n : a_1^A \ge b_{l-1}^B \land a_1^A \le a_l^B \land b_1^A \ge a_k^B \land b_1^A \le b_k^B \\
  \llbracket a_1^B, b_1^B, \ldots, a_l^B, b_k^B, a_{k+1}^B, \dots \rrbracket & \text{si } \exists l<k<n : a_1^A \ge a_l^B \land a_1^A \le b_l^B \land b_1^A \ge a_k^B \land b_1^A \le b_k^B \\
  \llbracket a_1^B, b_1^B, \ldots, a_l^B, b_1^A, a_{k+1}^B, \dots \rrbracket & \text{si } \exists l<k<n : a_1^A \ge a_l^B \land a_1^A \le b_l^B \land b_1^A > b_k^B \land b_1^A < a_{k+1}^B \\
  \llbracket a_1^B, b_1^B, \ldots, a_l^B, b_l^B, a_1^A, b_1^A, a_k^B, \dots \rrbracket & \text{si } \exists l<k<n : a_1^A > b_l^B \land a_1^A < a_{l+1}^B \land b_1^A > b_{k-1}^B \land b_1^A < a_k^B
\end{cases}
\end{align*}

Para el caso $m = 1$ o $m = 0$ y especiales queda demostrado el cerramiento de $B$ bajo esta suma reducida. 
El caso siguiente se construye con facilidad por recurrencia en cualquier número finito de pasos. 

Para cualquier $m > 1$:

\[
B + A := \begin{cases}
  \sum_{i=1}^{m} \left( B + \llbracket a_i^A, b_i^A \rrbracket \right) & \text{si } 1 \le i \le m, A_0 := \llbracket \varnothing \rrbracket, A_{i+1} := A_i + \llbracket a_i^A, b_i^A \rrbracket
\end{cases}
\]

Queda demostrado que toda suma da como resultado un conjunto finito de intervalos genéricos.

A su vez el producto lo vamos a definir de forma recursiva también, comenzando primero con $A$ siendo la clase de un solo intervalo genérico, o la clase del vacío, además de algunos casos especiales.

\[
B \cdot A := \begin{cases}
  \llbracket \varnothing \rrbracket & \text{si } A = \llbracket \varnothing \rrbracket \lor B = \llbracket \varnothing \rrbracket \\
  A & \text{si } B = 1 \\
  B & \text{si } A = 1 \\
  \llbracket \varnothing \rrbracket & \text{si } \exists A \in A, B \in B : B \subseteq \overline{A} \lor A \subseteq \overline{B} \\
  A & \text{si } \exists A \in A, B \in B : A \subseteq B \\
  B & \text{si } \exists A \in A, B \in B : B \subseteq A
\end{cases}
\]

Para el caso general de $m = 1$:

\begin{align*}
B \cdot A &:= \begin{cases}
  \llbracket a_1^A, a_1^B \dots \rrbracket & \text{si } a_1^A \le a_1^B \land \exists k<n : b_1^A \ge b_{k-1}^B \land b_1^A < a_k^B \\
  \llbracket a_1^B, b_1^B, \ldots, a_k^B, b_1^A \rrbracket & \text{si } a_1^A \le a_1^B \land \exists k<n : b_1^A \ge a_k^B \land b_1^A \le b_k^B \\
  \llbracket a_1^A, b_1^B, \ldots, a_{k-1}^B \dots \rrbracket & \text{si } a_1^A \ge a_1^B \land a_1^A \le b_1^B \land \exists k<n : b_1^A \ge b_{k-1}^B \land b_1^A < a_k^B \\
  \llbracket a_1^A, b_1^B, a_2^B, b_2^B, \ldots, a_k^B, b_1^A \rrbracket & \text{si } a_1^A \ge a_1^B \land a_1^A \le b_1^B \land \exists k<n : b_1^A \ge a_k^B \land b_1^A \le b_k^B \\
  \llbracket a_1^A, b_l^B, \dots, a_k^B, b_1^A \rrbracket & \text{si } \exists l<k<n : a_1^A \ge a_l^B \land a_1^A \le b_l^B \land b_1^A \ge a_k^B \land b_1^A \le b_k^B \\
  \llbracket a_1^A, b_l^B, \dots, a_{k-1}^B, b_{k-1}^B \rrbracket & \text{si } \exists l<k<n : a_1^A \ge a_l^B \land a_1^A \le b_l^B \land b_1^A > b_{k-1}^B \land b_1^A < a_k^B
\end{cases}
\end{align*}

Para el caso $m > 1$:

\[
B \cdot A := \begin{cases}
  \sum_{i=1}^{m} \left( B \cdot \llbracket a_i^A, b_i^A \rrbracket \right) & \text{si } 1 \le i \le m, A_0 := \llbracket \varnothing \rrbracket, A_{i+1} := A_i + \llbracket a_i^A, b_i^A \rrbracket
\end{cases}
\]
\end{enumerate}

Y queda demostrado que la forma del conjunto producto es una clase de
  unión finita de subintervalos genéricos. Luego pertenece a nuestro
  álgebra de Boole.

  Este sistema es muy parecido a un álgebra de conjuntos subálgebra de
  algún conjunto po\-tencia, por lo que es fácil determinar que se trata
  de un álgebra de Boole. Sin embargo su cardinal es
  \[\#{\left( {\left( {\mathbb{Q} \times \mathbb{Q}} \right) \times \left( {\mathbb{N} \times \mathbb{N}} \right)} \right) = \#}\left( \mathbb{N} \right)\].
  Veámoslo:

  \[{{\#\left( \mathbf{\mathrm{I}}_{\mathbb{Q}}^{\mathbf{\mathrm{GEN}}} \right)} = {\#\left( {\mathbf{\mathrm{\cup}}\begin{Bmatrix}
  {{⟦\varnothing ⟧},} & & & & & & \\
  {{⟦{a_{1,}b_{1}}⟧},} & & & & & & \\
  {{⟦{a_{1,}b_{1}}⟧},} & {{⟦{a_{2,}b_{2}}⟧},} & & & & & \\
  {{⟦{a_{1,}b_{1}}⟧},} & {{⟦{a_{2,}b_{2}}⟧},} & {{⟦{a_{3,}b_{3}}⟧},} & & & & \\
   \vdots & \vdots & \vdots & \ddots & & & \\
  {{⟦{a_{1,}b_{1}}⟧},} & {{⟦{a_{2,}b_{2}}⟧},} & {{⟦{a_{3,}b_{3}}⟧},} & \ldots & {{⟦{a_{n},b_{n}}⟧},} & & \\
   \vdots & \vdots & \vdots & \ddots & \vdots & \ddots & \\
  \ldots & \ldots & \ldots & \ldots & \ldots & \ldots & \ldots
  \end{Bmatrix}} \right)}} \leq\]\[{{\leq {1 + {\#\left( {\underset{i \in \mathbb{N}}{\mathbf{\mathrm{\cup}}}\underset{j \in \mathbb{N}}{\mathbf{\mathrm{\cup}}}{\mathbb{Q} \times \mathbb{Q}}} \right)}}} = \#}{\left( {{{\mathbb{N} \times \mathbb{N}} \times \mathbb{Q}} \times \mathbb{Q}} \right) = \#}{\mathbb{N} = \aleph_{0}}\]

  En definitiva es un álgebra numerable (del mismo cardinal que los
  números naturales). Puesto que el cardinal de los números naturales
  \[\mathbb{N}\]no es conmensurable con el de la potencia de ningún
  conjunto (es del cardinal infinito más pequeño posible y ningún
  conjunto finito tiene como potencia uno infinito), no existe ningún
  conjunto para el cual esta álgebra de Boole sea semejante (isomorfa) a
  un álgebra de las potencias de un conjunto. Este ejemplo será de
  utilidad más adelante, además de darnos un curioso ejemplo de álgebra
  de Boole nada común.

  \textbf{{[}Ejemplo 10{]}} El álgebra de Boole de los subconjuntos finitos y cofinitos de los números naturales ($\mathbb{N}$). Consideremos el conjunto base $B = \{ X \subseteq \mathbb{N} \mid X \text{ es finito o } \mathbb{N} \setminus X \text{ es finito} \}$. Las operaciones son las habituales de la teoría de conjuntos:
  \begin{itemize}
      \item \textbf{Suma ($+$):} La unión de conjuntos ($\cup$).
      \item \textbf{Producto ($\cdot$):} La intersección de conjuntos ($\cap$).
      \item \textbf{Complemento ($\overline{X}$):} El complementario relativo a $\mathbb{N}$ ($\mathbb{N} \setminus X$).
      \item \textbf{Cero ($0$):} El conjunto vacío ($\emptyset$).
      \item \textbf{Uno ($1$):} El conjunto de los naturales ($\mathbb{N}$).
  \end{itemize}
  Es inmediato comprobar que el complemento de un conjunto finito es cofinito (y viceversa), y que la unión o intersección de dos conjuntos finitos/cofinitos sigue produciendo un conjunto finito o cofinito. Por tanto, este conjunto es cerrado bajo las operaciones topológicas y constituye un álgebra de Boole de cardinalidad infinito numerable (igual que el Ejemplo 9).

  Sin embargo, es matemáticamente crucial señalar que \textbf{esta álgebra NO es isomorfa a la del Ejemplo 9}. Mientras que el Ejemplo 9 carece por completo de átomos (cualquier subintervalo racional puede subdividirse en dos más pequeños), esta álgebra del Ejemplo 10 \textbf{sí es atómica}: sus átomos son precisamente los conjuntos compuestos por un único número natural, $\{n\}$. Esta genialidad geométrica demuestra que, en el infinito, pueden existir álgebras de Boole del mismo cardinal que son estructuralmente distintas (algo que, como demostramos en el Capítulo 7, es imposible en las álgebras finitas).

  \textbf{{[}Ejemplo 11{]}} El álgebra de las funciones de un álgebra de
  Boole sobre otra. Supongamos \[f:{B\rightarrow B}'\]dónde \[f\] es una
  función. Llamaremos
  \[{\mathtt{F}{({B,B'})}} = {\{{f:{B\rightarrow B}' \mid \forall{x \in B}\exists!{y \in B}'f{{(x)} = y}}\}}\]
  a nuestro conjunto de Boole. \[f_{0'}\] es la función que asigna el
  cero de \[B'\] a cualquier elemento de \[B\]. Igualmente \[f_{1'}\] es
  la función que asigna el uno de \[B'\] a cualquier elemento de \[B\].
  Las operaciones internas a introducir son:

  La suma de funciones:
  \[\forall{x \in B}{\lbrack{f + g}\rbrack}{(x)}{: = f}{{(x)} + g}{(x)}\]
  .

  La multiplicación de funciones:
  \[\forall{x \in B}{\lbrack{f \cdot g}\rbrack}{(x)}{: = f}{{(x)} \cdot g}{(x)}\].

  La función complemento:
  \[\forall{f \in \mathtt{F}}{({B,B'})}\forall{x \in B}\overline{f}{(x)}{: = \overline{f{(x)}}}\].

  De esta álgebra podemos entresacar otros conjuntos de funciones
  interesantes, como. Por ejemplo:

  \[{\mathtt{\mathit{Hom}}{({B,B'})}} = \begin{Bmatrix}
  f & : & B & \rightarrow & {B'} & \mathbf{\mid} \\
  {\forall{x \in B}} & {\exists!} & {{y \in B}'} & {f{{(x)} = y}} & \mathbf{\mathrm{:}} & \\
   & & & & {f{{(0)} = 0}'} & \land \\
   & & & & {f{{(1)} = 1}'} & \land \\
   & & {\forall x,{y \in B}} & & {f{{({x + y})} = f}{{(x)} + f}{(y)}} & \land \\
   & & {\forall x,{y \in B}} & & {f{{({x \cdot y})} = f}{{(x)} \cdot f}{(y)}} & \land \\
   & & {\forall{x \in B}} & & {f{{(\overline{x})} = \overline{f{(x)}}}} & 
  \end{Bmatrix}\]

  y aún otros subconjuntos más pequeños serían las inyecciones de los
  anteriores homomorfismos. Si \[{B'} \equiv B\]entonces uno de los
  conjuntos de aplicaciones más interesantes son los endomorfimos o
  isomorfimos en sí mismo.

  Además el kernel de cualquier homomorfismo es un subálgebra de \[B\].
  Los homomorfismos de las álgebras booleanas tienen propiedades
  interesantes que no veremos aquí.
